% CLAUDE REPORT::B3-RELATIVIZE-0001::2026-11-12T02:53Z
% !TEX program = xelatex
% Compiles with XeLaTeX or pdfLaTeX (ASCII source).
\documentclass[11pt]{article}
\usepackage[margin=1in]{geometry}
\usepackage{amsmath,amssymb,amsthm}
\usepackage[hidelinks]{hyperref}
\newtheorem{theorem}{Theorem}
\newtheorem{lemma}[theorem]{Lemma}
\newtheorem{corollary}[theorem]{Corollary}
\newtheorem{remark}[theorem]{Remark}
\newcommand{\SAT}{\mathrm{SAT}}
\newcommand{\Wb}{\overline W}
\newcommand{\TQBF}{\mathrm{TQBF}}
\newcommand{\Rmix}{R_{\mathrm{mix341}}}
\title{The B3 Bridge Has No Relativizing or Algebrizing Proof\\
unless the Fixed Calculus Is Polynomially Bounded}
\author{Claude (scheduled worst-case-derivation run)}
\date{11 November 2026, 18:53 PST}
\begin{document}
\maketitle

\section*{Status}
\textbf{P versus NP remains OPEN.} Theorems~1--3 and Corollaries~4--6 below are
\textbf{PROVED/EXACT} from classical facts (Cook--Reckhow 1979; Stockmeyer/Meyer
PSPACE-completeness of TQBF; Baker--Gill--Solovay 1975; Aaronson--Wigderson 2009).
They are barrier statements about the \emph{shape} of any proof of the bridge B3; they
prove no lower bound and no bridge. B1, B2, B3 for $\Rmix$ and a superpolynomial
$\Wb_{\Rmix}$ remain \textbf{OPEN}.

\section*{Shared objective and the question addressed}
For a fixed sound, complete, effectively checkable calculus $R$ with unit (or
nonnegative) node weights,
\[
 \Wb_R(N)=\max_{F\in\mathrm{UNSAT},\ |\langle F\rangle|\le N}\ \min_{\pi:F\vdash_R\bot}\operatorname{cost}_R(\pi).
\]
The inner minimum grants free knowledge of the optimal reduction sequence; search cost is
not charged. The primary bridge is
\[
 \mathrm{B3}_R:\qquad \SAT\in \mathrm P\ \Longrightarrow\ \exists C,k\ \forall N:\ \Wb_R(N)\le C(N+1)^k .
\]
\emph{Question.} Which proof techniques can possibly establish $\mathrm{B3}_R$ for the
project calculus $\Rmix$? LSPV-1/548 (main.tex, Resolution oracle note) proved that
B3 fails relative to TQBF for ordinary Resolution and that, \emph{conditional on} a future
superpolynomial lower bound for $\Rmix$, any B3 proof must be nonrelativizing. This run
removes that condition and adds the algebrization version.

\section*{Falsifiers proposed before work}
\begin{enumerate}
\item[F1] $\Wb_R$ might secretly depend on the oracle (then the TQBF instantiation would not
collapse the consequent). \emph{Ruled out:} $R$, $F$ and $\operatorname{cost}_R$ are oracle-free.
\item[F2] Node cost of $\Rmix$ might be exponentially smaller than encoded length, so
``$\Wb$ polynomial'' would not give a polynomially bounded Cook--Reckhow system.
\emph{Ruled out by Lemma~2.}
\item[F3] ``SAT $\in$ P$^O$'' might be read as SAT$^O$ (oracle-circuit SAT); then
TQBF would not obviously satisfy it. \emph{Addressed:} the antecedent concerns ordinary
CNF-SAT, the object of $\Wb_R$; the oracle-SAT reading is treated in Remark~7.
\item[F4] The algebrization version might need $P^{\tilde A}\supseteq \mathrm{NP}^A$
rather than ordinary SAT. \emph{Addressed in Theorem~3}: only ordinary SAT $\in P^{\tilde A}$ is
needed, and it holds because $\tilde A$ agrees with $A$ on Boolean points.
\end{enumerate}

\section{Relativized bridges}
For an oracle $O\subseteq\{0,1\}^*$ write $H^O$ for: \emph{there is a deterministic
polynomial-time oracle machine $A^O$ deciding ordinary encoded CNF-SAT (rejecting
malformed inputs)}. The relativized bridges are
\[
 \mathrm{B3}_R^O: H^O\Rightarrow \Wb_R\le\mathrm{poly},\qquad
 \mathrm{B2}_R^O: \forall A^O\text{ deciding SAT}\ \exists p\ \forall N:\ \Wb_R(N)\le p(N+T^{\mathrm{wc}}_{A^O}(N)),
\]
and $\mathrm{B1}_R^O$ is the pointwise analogue. A proof technique \emph{relativizes} if
whatever it proves for B$i_R$ it proves for B$i_R^O$ for every $O$ (Baker--Gill--Solovay).

\begin{theorem}[TQBF collapse of the relativized bridge]
For every oracle-free calculus $R$: $\mathrm{B3}_R^{\TQBF}$, $\mathrm{B2}_R^{\TQBF}$ and
$\mathrm{B1}_R^{\TQBF}$ are each equivalent to the oracle-free statement
``$\Wb_R(N)\le\mathrm{poly}(N)$''.
\end{theorem}
\begin{proof}
$H^{\TQBF}$ holds: a machine parses $F$ in polynomial time (rejecting malformed input),
writes the closed QBF $\exists x_1\cdots\exists x_n F$ and queries the oracle once; total
worst-case time is polynomial. Hence $\mathrm{B3}^{\TQBF}_R$ reduces to its consequent.
For B2/B1, this one decider has $T^{\mathrm{wc}}(N)\le\mathrm{poly}(N)$, so the conclusion
for it gives $\Wb_R\le\mathrm{poly}$; conversely, if $\Wb_R$ is polynomial then the B2/B1
inequalities hold for every decider since $T\ge 0$. The consequent mentions no oracle
because $R$, the formulas and the cost function are oracle-free.
\end{proof}

\begin{lemma}[Node cost controls encoded length in $\Rmix$]
Let $\pi$ be an $\Rmix$ refutation of $F$, $|\langle F\rangle|=N$, with $q$ charged nodes.
Then $\pi$ has an encoding of $O\bigl(q\,(N+q)\log(N+q)\bigr)$ bits, checkable in time
polynomial in $N+q$.
\end{lemma}
\begin{proof}
$F$ has at most $N$ variables; each DEF node adds one variable, so at most $N+q$ variables
occur. Every line is a clause, a GF(2) equation (variable mask plus right-hand side), or a
truth table on at most three variables, hence has $O((N+q)\log(N+q))$ bits; parent pointers
cost $O(\log q)$. Each rule (RES, AND3 with 8-row check, XOR, LOCAL-CONVERT with constant
support, DEF) is checked in time polynomial in its line sizes.
\end{proof}

\begin{theorem}[Relativizing B3 proofs prove NP $=$ coNP]
If a relativizing argument proves $\mathrm{B3}_{\Rmix}$ (or B2 or B1), then the same
argument proves unconditionally that $\Rmix$ is a polynomially bounded Cook--Reckhow
proof system, hence NP $=$ coNP.
\end{theorem}
\begin{proof}
A relativizing argument yields $\mathrm{B3}^{\TQBF}_{\Rmix}$, so by Theorem~1,
$\Wb_{\Rmix}(N)\le C(N+1)^k$. By Lemma~2 every UNSAT $F$ has an $\Rmix$ refutation of
length polynomial in $N$, verifiable in polynomial time; $\Rmix$ is sound (each rule is
sound and DEF is conservative) and complete (RES). Thus UNSAT $\in$ NP, and since UNSAT is
coNP-complete, NP $=$ coNP (Cook--Reckhow).
\end{proof}

\begin{theorem}[Algebrization version]
Let $A=\TQBF$ and let $\tilde A$ be any extension of $A$ to a finite field
(Aaronson--Wigderson sense: $\tilde A$ agrees with $A$ on Boolean points). Then
ordinary SAT $\in P^{\tilde A}$. Hence any argument establishing the algebrized bridge
``SAT $\in P^{\tilde A}\Rightarrow \Wb_R\le\mathrm{poly}$'' for all such pairs proves
$\Wb_R\le\mathrm{poly}$ unconditionally; for $R=\Rmix$ it proves NP $=$ coNP.
\end{theorem}
\begin{proof}
The decider of Theorem~1 queries only a Boolean point (the encoding of a closed QBF), where
$\tilde A$ returns $A$'s value. The rest is Theorem~1 and Theorem~2's argument.
\end{proof}

\begin{corollary}[Calculi with known lower bounds]
If $\Wb_R$ is unconditionally superpolynomial (Resolution: Haken 1985; also cutting planes,
regular Res$(\oplus)$, bounded-depth Frege), then $\mathrm B3_R$ is \emph{equivalent} to
$\SAT\notin\mathrm P$ (LSPV-1/398) and admits no relativizing or algebrizing proof.
\end{corollary}

\begin{corollary}[Where the nonrelativizing content must sit]
In the two-part strategy ``superpolynomial $\Wb_R$'' plus ``$\mathrm B3_R$'', the
lower-bound half is an oracle-free statement about a fixed calculus and is not touched by
these barriers. Unless $\Rmix$ is polynomially bounded, \emph{all} nonrelativizing and
nonalgebrizing content must be carried by the bridge.
\end{corollary}

\begin{corollary}[Reflection route is consistent with the barrier]
The known positive route (Cook 1975; Krajíček--Pudlák 1989): if a polynomial-time SAT
decider $A$ has a PV-provable (equivalently short EF-provable) UNSAT-soundness statement,
then $\Rmix\supseteq\mathrm{ER}$ p-simulates the system ``run $A$'', and B3 follows. This
route does \emph{not} relativize: for $A^{\TQBF}$, the soundness statement contains oracle
answers that $\Rmix$ cannot certify. So the barrier does not exclude it; it explains why
it is the only kind of route presently on file. Its missing premise --- that SAT $\in$ P
implies PV-provable correctness of \emph{some} polynomial-time decider --- is \textbf{OPEN},
and the web search below found no theorem deriving EF boundedness from $\mathrm P=\mathrm{NP}$.
\end{corollary}

\begin{remark}[Oracle-SAT reading]
If the antecedent is relativized to SAT$^O$ (CNFs with oracle gates), the consequent would
need an oracle calculus $R^O$; the theorems above concern the project's oracle-free $\Rmix$
and ordinary CNF, which is the target of $\Wb_R$. No claim is made for the other reading.
\end{remark}

\section*{Literature check (this run)}
\begin{itemize}
\item Query ``lower bounds resolution over parities DAG-like 2025'': no superpolynomial
lower bound is known for general DAG-like Res$(\oplus)$; exponential bounds exist for regular
Res$(\oplus)$ (Efremenko--Garl\'ik--Itsykson, STOC 2024), bounded-depth Res$(\oplus)$
(Alekseev--Itsykson 2025; Itsykson--Knop; Byramji--Impagliazzo for bit-PHP, ECCC 2025).
Relevance: the DEF-free fragment of $\Rmix$ is p-simulated by Res$(\oplus)$ (clauses and
single equations are linear clauses; AND3 tables are constant-size clause sets), so a
lower bound for that fragment would sit inside an open problem; with DEF, $\Rmix\ge$ ER and
a lower bound would be an Extended-Frege-strength result. Sources:
arXiv 2507.23008, arXiv 2511.20023, ECCC TR25-039/116/118/188 (abstracts read).
\item Query ``P = NP implies Extended Frege polynomially bounded reflection'': found
Cook--Reckhow (p-bounded system iff NP $=$ coNP), Krajíček--Pudlák optimal-system results
(NE $=$ coNE $\Rightarrow$ optimal systems), EF+Ref$_Q$ simulates $Q$, and P $=$ NP implies
automatability of every system. No source derives EF (or $\Rmix$) boundedness from
P $=$ NP; consistent with Corollary~6. Failure to find is not evidence of nonexistence.
\end{itemize}

\section*{CORPUS REVIEWED}
AGENTS.md (all addenda including the 2026-11-09 knot/link handoff pointer);
research/coordination/2026-11-02\_1204\_WORST\_CASE\_MINIMUM\_DERIVATION\_FIVE\_ROLE\_REFOCUS.md
(shared $\Wb_R$ directive, B1/B2/B3 definitions, LSPV-2/341 calculus);
LSPV-2/341 report (exact rules of $\Rmix$); LSPV-1/398 (Resolution B3 equivalence);
LSPV-1/355, /361, /367, /368, /372 (nonrelativizing-B3 remarks); LSPV-1/548 prewrite
backup (TQBF Resolution oracle note); Codex 2026-11-11 1203 and 1331 reports (status tables,
reflection/witnessing OPEN, Pich--Santhanam 2023); root main.tex highest-value markers;
prior Claude reports in research/agents/claude/. Repository search for
``relativi''+``B3'', ``Baker'', ``algebriz'', ``Res($\oplus$)''.

\section*{DUPLICATION / CONTRADICTION AUDIT}
\emph{Duplicate parts:} the Resolution TQBF counterexample and the conditional
nonrelativization remark are LSPV-1/548; the Resolution B3 $\equiv$ SAT $\notin$ P
equivalence is LSPV-1/398; the reflection route is Cook 1975 / KP89 as recorded in
LSPV-1/357 and Codex 1331. \emph{New here:} (i) removal of the lower-bound condition:
a relativizing B3 proof for $\Rmix$ would itself prove NP $=$ coNP (Theorem~2, via
Lemma~2); (ii) the algebrizing version with an explicit Boolean-point argument
(Theorem~3); (iii) the explicit allocation statement (Corollary~5). No contradiction with
the corpus was found.

\section*{WHAT CHANGED RELATIVE TO ROOT MAIN}
Adds an unconditional technique barrier on B3 for the unchanged $\Rmix$. No status in
root main.tex changes: B3 and the superpolynomial lower bound stay OPEN. This Claude run
owns no main.tex region and does not edit it; LSPV-5 may adopt it at synthesis.

\section*{ADOPTED / REJECTED / RETIRED PRIOR COMPONENTS}
Adopted: $\Rmix$ as specified in LSPV-2/341; LSPV-1/398 and /548 results. Rejected: none.
Retired: the condition ``if $\Wb_{\Rmix}$ proves superpolynomial'' in the /548 remark is no
longer needed for the nonrelativization conclusion (it is now implied unless NP $=$ coNP).

\section*{STALE OR OVERSTATED CLAIMS FOUND}
None new. Reminder carried forward: Theorem~2 does not say B3 is false or unprovable; it
says a relativizing proof would prove the (widely disbelieved, open) statement NP $=$ coNP.

\section*{Post-quantum technique ledger}
Item 5 (worst-case/average-case, black-box extraction): a black-box reduction that turns
any SAT decider into short $\Rmix$ proofs would typically relativize, so by Theorem~2 it
cannot exist unless NP $=$ coNP. Classification: FAILED TRANSFER for black-box
extraction; reflection-based (white-box, provability) routes remain CANDIDATE.

\section*{Status table}
\begin{tabular}{ll}
TQBF collapse of relativized B1/B2/B3 (Thm 1) & PROVED/EXACT\\
$\Rmix$ node cost vs length (Lemma 2) & PROVED/EXACT\\
Relativizing B3 for $\Rmix$ $\Rightarrow$ NP $=$ coNP (Thm 2) & PROVED/EXACT\\
Algebrizing version (Thm 3) & PROVED/EXACT\\
B1/B2/B3 for $\Rmix$ & OPEN\\
Superpolynomial $\Wb_{\Rmix}$ & OPEN (EF-strength)\\
P versus NP & OPEN
\end{tabular}

\section*{Single next target}
Find a statement strictly weaker than ``PV proves correctness of a polynomial-time SAT
decider'' that (a) is implied by SAT $\in$ P, (b) fails relative to TQBF, and (c) still
gives short $\Rmix$ proofs of a universal soundness family. Any candidate must say
which nonrelativizing ingredient (arithmetization, provability, explicit code inspection)
it uses.

\section*{References}
S.~Cook and R.~Reckhow, J.\ Symbolic Logic 44 (1979). S.~Cook, STOC 1975.
J.~Krajíček and P.~Pudlák, J.\ Symbolic Logic 54 (1989). T.~Baker, J.~Gill, R.~Solovay,
SIAM J.\ Comput.\ 4 (1975). S.~Aaronson and A.~Wigderson, ACM TOCT 1 (2009).
A.~Haken, TCS 39 (1985). K.~Efremenko, M.~Garl\'ik, D.~Itsykson, STOC 2024.
Y.~Alekseev, D.~Itsykson, STOC 2025. ECCC TR25-039, TR25-116, TR25-118, TR25-188;
arXiv 2507.23008, 2511.20023.

\end{document}
% END_OF_VERIFIED_REPORT::CLAUDE::B3-RELATIVIZE-0001
